% frontmatter/como-usar.tex -- instrucciones para el adulto.
\section*{Cómo usar este cuaderno}

Este cuaderno es para una niña o un niño de cinco o seis años que
empieza con los números: reconocerlos, contarlos y escribirlos -- y, a lo
largo del año, compararlos, sumar y restar. Es de la misma familia que
\emph{Aprendo a leer}: los mismos personajes (Lucía, su hermano Dani, su
perro Toby, Mamá, Papá y la abuela Rosa) y los mismos temas semana a
semana, así que quien lleve los dos cuadernos cuenta esa semana lo que
lee.

Hay una página para cada día laborable del año, de lunes a viernes,
también en vacaciones. Cinco o diez minutos al día son suficientes.
Cada página sigue siempre la misma estructura:

\begin{itemize}[leftmargin=6mm]
\item \textbf{Fecha}, \textbf{Hecho} y \textbf{Notas}: para el adulto.
  Marcar cada día si lo ha hecho sola, con ayuda o con dificultad no
  cuesta nada, y al cabo del año es un registro de todo lo que ha
  avanzado.
\item \textbf{El número de hoy}, en el recuadro verde: el número grande,
  su nombre, el \textbf{marco de diez} -- dos filas de cinco casillas,
  con un punto por cada uno: así se ve de un vistazo cuánto es un
  número (el 5 es una fila entera; el 7, una fila y dos más) -- y esas
  cosas, dibujadas. Debajo, una frase de la historia, que lee el adulto:
  el número de hoy, en casa de Lucía.
\item Una \textbf{actividad}, que casi siempre termina coloreando,
  rodeando o dibujando. Las instrucciones (en letra pequeña y gris)
  \textbf{las lee un adulto} en voz alta.
\end{itemize}

\subsection*{Las actividades}

{\small
\begin{itemize}[leftmargin=6mm, itemsep=1pt, topsep=2pt]
\item \textbf{\lblTraza} (los lunes, cuando llega un número nuevo):
  repasar el número con el dedo por encima de los puntos, luego con un
  lápiz, y después escribirlo solo en la línea. Los puntos son el
  contorno del número de verdad, el mismo que se lee en el resto del
  cuaderno.
\item \textbf{\lblColorea}: colorear exactamente las que dice, ni una
  más -- las demás se quedan en blanco.
\item \textbf{\lblRodea}: entre varios grupos de cosas, rodear el que
  tiene las que dice. Los grupos están puestos como los puntos de un
  dado, para que se vean de un vistazo.
\item \textbf{\lblCuenta}: contar señalando cada cosa con el dedo, una
  sola vez cada una, y rodear el número.
\item \textbf{\lblBusca}: rodear el número cada vez que aparece, entre
  otros números y formas.
\item \textbf{\lblDibuja}: dibujar las cosas que dice, ni más ni menos.
\item \textbf{\lblRepasa} (los viernes): lo que ya sabe hacer esa
  semana, para marcarlo, y un dibujo. Cada diez páginas hay un pequeño
  cartel de celebración.
\end{itemize}}

\subsection*{Cómo ayudar}

{\small
\begin{itemize}[leftmargin=6mm, itemsep=1pt, topsep=2pt]
\item Contar es tocar: que señale cada cosa con el dedo al contarla, y
  una sola vez. Si se salta una o cuenta dos veces la misma, contad
  juntos, despacio.
\item Di el nombre del número cada vez que lo traza o lo encuentra: el
  número se ve, se dice y se escribe.
\item Si algo le cuesta, se puede volver a hacer al día siguiente con
  cosas de casa: cucharas, botones, zapatos.
\item Al final de cada trimestre hay una medalla, y al final del
  cuaderno, un diploma y una \textbf{clave de respuestas} para las
  páginas de \lblRodea, \lblCuenta\ y \lblBusca.
\end{itemize}}
\newpage
